%  ADDED 2026-07-06: full subsection "The frame" (sec:frame), between
%  the two-sphere representation and the spin ladder: the frame defined
%  (oriented coordinatisation; map-grid not territory), gauge-per-sphere
%  / constraint-per-pair, discreteness as the promotion, the electroweak
%  inventory frame-borne, the lepton/quark certification asymmetry,
%  turned-never-held with the budget theorem (sec:budget-converse), the
%  parity consonance (flagged; op:chirality), the bundle/cocycle reading
%  (flagged; travels with op:filatov-extension). "Frame" is now a
%  defined term citable via sec:frame.
%
%  RE-AUDITED 2026-07-06 (final pass, tenth stop, light leg): position
%  10 adjacency claims TRUE (both "previous section" phrases refer to
%  entanglement-classes content -- CK composition, class reading);
%  frame-subsection seams verified ("previous subsection" = the
%  two-sphere representation, "next subsection" = the spin ladder,
%  both correct); op:chirality and op:filatov-extension confirmed
%  DEFINED in the particle zoo -- the frame's references land; all ref
%  targets exist; no textquote; header Goedel umlaut normalised, file
%  now pure ASCII. No prose changes required: the deep audit and the
%  frame insertion stand.
% =====================================================================
%  Section: A Fermion is Three Qubits   (Plan Stage 1)
%  Paper 1, "It from Bit via G\"odel"
%  Drafted 2026-06-08. Follows \S sec:entanglement-classes.
%
%  \input fragment. Assumes biblatex/\autocite, amsthm `proposition'.
%  Bib resources: hopf_particle_physics_refs_2026-06-08.bib,
%                 entanglement_classes_refs_2026-06-08.bib,
%                 register_refs_2026-06-08.bib (new).
%
%  Cross-reference labels (placeholders -- re-point before compiling):
%    sec:hopf-justification     -- Hopf fibrations / Adams termination / prop:gsm
%    sec:entanglement-classes   -- GHZ/W theorems, "split on the first qubit"
%    prop:gsm, prop:termination -- from the Hopf-justification section
%    sec:fanout-direction       -- the computability split / measurement
%    sec:charges                -- Stage 2 (reading the quantum numbers), forward
%    sec:generations            -- Stage 4 (three generations), forward
%    sec:masses                 -- Stage 5+ (mass), forward
%
%  Terminology follows terminology_reconciliation_entanglement_dictionary_2026-06-08:
%  "Bell" is the entanglement class, never a synonym for SU(2); the $\mathbb{H}$ rung is
%  named as such. Dual-audience convention (QC view / GT view / Bridge) per the plan.
% =====================================================================

\section{A Fermion is Three Qubits}
\label{sec:register}

\subsection{Particles are channels, not objects}
\label{sec:register-preface}

A caveat first, because it governs this section and every one that follows. The language we are
about to adopt --- the ``internal structure'' of a fermion, a ``three-qubit register'', a qubit
that ``carries'' colour --- is convenient and we will use it throughout, but it must not be read
literally. There is no little ball, existing in its own right, with hypercharge and colour painted
on its surface. That image is a classical inheritance: in everyday physics we deal with objects
that persist independently of their interactions and bear properties as attributes, and we carry
the habit into quantum physics, where it quietly misleads.

In this framework a particle is not an object but a \emph{channel of communication between events}.
The fundamental layer is the self-referential network of events of \S\ref{sec:self-reference}; a
particle is a connection across that network, and what we call its properties are features of what
the channel can carry and of how it composes with other channels, not marks on a thing. The
``particle'' is the interpretation our classically-trained intuition places on the information the
channel conveys --- a useful repackaging, not the discovery of a blob.

This is not a philosophical ornament bolted onto the mathematics; it is what makes the mathematics
of the previous section the \emph{right} mathematics. The compositional structure of Coecke and
Kissinger \autocite{coeckekissinger2010compositional,CoeckeKissinger2017} is a calculus of processes: systems are
wires, interactions are boxes, and the content of the theory is how they compose. A particle, in
that calculus, is a wire carrying information between the events it connects, and the GHZ and W
Frobenius structures are statements about how such a wire may branch, copy and merge --- properties
of a channel, not labels on an object. When \S\ref{sec:entanglement-classes} let the entanglement
\emph{class} decide whether something is a quark or a lepton, it was already committed to this
reading: a quark is not a coloured ball but a channel whose composition behaviour is of the GHZ
type.

Both of our audiences are, in fact, already thinking this way, whatever the popular picture
suggests. The quantum-computational reader works with states and gates --- preparations and
processes; the gauge-theoretic reader works with fields and vertices --- again processes, and the
events they connect. It is only the pop-physics image of a particle as a tiny labelled sphere that
is the outlier, and it is exactly that image we are setting aside.

So the phrases that follow are bookkeeping. ``The three-qubit register'' is a way of organising the
information a channel can carry, written in a language --- qubits, Bloch spheres, gates --- that
maps cleanly onto both audiences' habits and onto our shared classical experience. The three qubits
are not three sub-things lodged inside a fourth thing. With that understood the convention is
genuinely useful and we adopt it without further apology, but the reader should mentally supply
``the channel we conventionally describe as'' in front of every ``the fermion'' below.

\subsection{What this section does}

The previous section read the entanglement \emph{class} of a three-qubit state as the
quark--lepton division. To go further --- to the quantum numbers of \S\ref{sec:charges} and the
masses of \S\ref{sec:masses} --- we need to make explicit the structure of the channel we
conventionally describe as a single fermion, written, by the caveat just stated, as a three-qubit
register with each qubit tied to one gauge sector. This section assembles that
register. It is worth saying plainly that the result is \emph{structural}: every ingredient (the
division-algebra tower, the Hopf fibrations, the two-sphere representation, the termination at
$\mathbb{O}$) is established in \S\ref{sec:division-algebras} and
\S\ref{sec:hopf-justification}, but the identification of the assembled register with a Standard
Model fermion is an argument from structure rather than a theorem, and we flag the assumptions as
they arise.

The section also fixes the dual reading the rest of the development relies on. The same object is,
for a reader coming from quantum computation, a specific three-qubit entangled state; and for a
reader coming from gauge theory, a particle carrying hypercharge, weak isospin and colour. Holding
both readings at once is the whole point, so each step below states a \textbf{QC view}, a
\textbf{GT view}, and the \textbf{Bridge} that is the single object both are describing.

\subsection{The three qubits and the Cayley--Dickson tower}

The channel we describe as a single fermion is organised, by the convention of
\S\ref{sec:register-preface}, into three qubits $A \otimes B \otimes C$, whose
internal algebras are the three non-trivial rungs of the Cayley--Dickson ladder of
\S\ref{sec:division-algebras}:
$\mathbb{C}$, $\mathbb{H}$, $\mathbb{O}$. Each rung is the previous one doubled, and each doubling
sacrifices a piece of algebraic structure: $\mathbb{C}$ is commutative and associative;
$\mathbb{H}$ is non-commutative but associative; $\mathbb{O}$ is neither.

\emph{QC view.} Three qubits, each an ordinary two-state system you already know how to write down
and rotate. \emph{GT view.} Three internal labels --- one abelian charge and two non-abelian ones.
\emph{Bridge.} The algebra a qubit is valued in fixes the gauge content it can carry; the tower is
the list of what structure survives at each level.

That there are exactly three such qubits is not a modelling choice. The same classification that
closed the gauge group in \S\ref{sec:hopf-justification} closes the register: there are exactly
three non-trivial normed division algebras above $\mathbb{R}$, exactly three nontrivial Hopf
fibrations, and the tower terminates at $\mathbb{O}$ (Proposition~\ref{prop:termination}). The
register inherits its length from the same theorem that forbids a fourth force.

\subsection{The qubits are nested, not side by side}

It is tempting to picture $A$, $B$, $C$ as three independent boxes. They are not. They are the
three levels of the \emph{iterated} octonionic Hopf fibration of the single three-qubit state, the
same nesting drawn by Szangolies \autocite{szangolies2025standardmodel} and exploited in the
``split on the first qubit'' of \S\ref{sec:entanglement-classes}.

\begin{proposition}[The register as an iterated fibration]\label{prop:iterated-hopf}
The unit sphere $S^{15}$ of the three-qubit state space is the total space of the octonionic Hopf
fibration $S^{15} \xrightarrow{S^7} S^8$; the $S^7$ fibre is itself the total space of
$S^7 \xrightarrow{S^3} S^4$; and the $S^3$ fibre is in turn $S^3 \xrightarrow{S^1} S^2$. The three
Hopf fibres are therefore nested, $S^1 \subset S^3 \subset S^7$, and they carry, respectively, the
groups $U(1)$, $SU(2)$, and --- through the $G_2$ stabiliser of a complex direction
(Proposition~\ref{prop:gsm}) --- $SU(3)$.
\end{proposition}

This is why peeling one qubit off the register drops one rung of the tower, exactly as in the
entanglement-classes construction. The three gauge groups are not bolted on; they are the three
fibres of one geometric object.

\subsection{One Bloch sphere per gauge sector}

Collecting the assignment gives the table that the rest of the development reads off. This table is
the deliverable of the stage: a reader of either background should be able to point at a row and
name both the qubit and the force.

\begin{center}
\begin{tabular}{lllll}
\hline
Qubit & CD rung & Hopf fibre & Gauge group & Internal charge \\
\hline
$A$ & $\mathbb{C}$ (commutative) & $S^1$ & $U(1)$ & hypercharge / electric \\
$B$ & $\mathbb{H}$ (non-commutative) & $S^3$ & $SU(2)$ & weak isospin \\
$C$ & $\mathbb{O}$ (non-associative) & $S^7$ & $SU(3)$ & colour \\
\hline
\end{tabular}
\end{center}

The ordering is fixed, not chosen: it is the order in which the Cayley--Dickson construction
sheds structure, matched to the order in which the gauge groups gain it. An abelian group sits on
the commutative rung ($U(1)$ on $\mathbb{C}$); a non-abelian group with an associative Lie algebra
sits on the non-commutative-but-associative rung ($SU(2)$ on $\mathbb{H}$), where left and right
quaternionic multiplication supply the two candidate actions and bipartite (Bell-class)
entanglement is naturally accommodated; and the group whose octonionic realisation is constrained
by non-associativity sits on the non-associative rung ($SU(3)$ on $\mathbb{O}$, as the $G_2$
stabiliser). We flag this as a \emph{structural matching} rather than a uniqueness theorem: it is
the natural assignment, and we know of no competitor, but ``natural'' is doing real work in that
sentence.

Two of the lost symmetries are worth naming now, because they are cashed out later, and the
division of labour between them is exact. Losing commutativity at the $\mathbb{H}$ rung supplies
the \emph{commutator}: once the order of two actions matters, the difference between the two
orderings is itself an action, and a closed algebra of such differences is precisely a non-abelian
gauge sector. The $W^{\pm}$ are that structure made physical --- isospin raising and lowering whose
compositions in opposite orders disagree, their commutator generating the third direction of
$SU(2)$. Losing associativity at the $\mathbb{O}$ rung supplies the \emph{associator} --- the
re-bracketing of composite processes, a capacity the quaternions do not possess --- and with it
chirality (triality breaking, \S\ref{sec:generations}) and the currency of mass
(\S\ref{sec:masses}). Re-ordering is the $\mathbb{H}$ deficit; re-bracketing is the $\mathbb{O}$
deficit; the two are distinct capacities, and this paper reserves the words accordingly.

One apparent paradox is worth defusing here, before the mass sections make the resolution precise.
If mass is the associator's currency, why are the $SU(2)$ bosons massive while the gluons ---
natives of the non-associative rung --- are massless? Because mass is not a deficit of a boson's
home rung: it is associator \emph{debt}, incurred relative to the selected complex direction
$\mathbb{C} \subset \mathbb{O}$ by whatever action fails to preserve that direction
(\S\ref{sec:masses}, \S\ref{sec:associator-higgs}). The $W^{\pm}$ are quaternionic in their
group law, but their action rotates the selection, and the obstruction to that rotation ---
measured in the octonionic ambient where the selection lives --- is their mass. The photon's action
preserves the selection and owes nothing. The gluons act within the colour sector, orthogonal to a
selection that is a colour singlet, and owe nothing either, for all their non-associative
citizenship. (The full classification of which bosons move the selection, photon and $Z^{0}$
included, is the business of \S\ref{sec:boson-masses}.) The register is thus not a set of labels
attached to an object: each algebraic deficit is a physical capacity of the channel --- and mass is
a debt against the selection, not a property of a rung.

\subsection{Drawing the register: the two-sphere representation}

The visualisation tool is the two-Bloch-sphere representation of Filatov and Auzinsh
\autocite{filatov2024towards}. They showed that a pure two-qubit entangled state can be drawn as a
pair of Bloch spheres, with the constraint that the coordinate axes of one sphere are
right-handed and those of the other left-handed. This opposite-handedness --- the \emph{Filatov
constraint} --- is a topological fact about the quaternionic Hopf fibration, not a drawing
convention.

\emph{QC view.} A multi-sphere picture of entanglement: each qubit gets a sphere, and entangled
pairs are tied together by opposite handedness. \emph{GT view.} The handedness is a physical
chirality carried by the internal state. \emph{Bridge.} The handedness relation between two spheres
is the entanglement between two gauge sectors.

For a fermion we draw three such spheres, one per row of the table. When two of the qubits are
entangled, their spheres carry opposite handedness. The case where all three pairs are
mutually entangled --- the W class --- demands opposite handedness on every edge of the triangle,
which the complete graph $K_3$ cannot consistently support: $K_3$ is not two-colourable. That
frustration is the seed of chirality and of the generation count, but it belongs to
\S\ref{sec:generations}; here we only need the representation itself and the fact that the
constraint is forced rather than imposed. One disclosure travels with the picture: Filatov's
theorem is for a single entangled pair, and the extension to three mutually constrained spheres ---
each pair carrying its own quaternionic Hopf structure --- is assumed here, its formal proof open;
the assumption is carried explicitly into \S\ref{sec:generations}.

\subsection{The frame}\label{sec:frame}

The sections that follow, and several that have gone before, lean on a word --- \emph{frame} ---
that deserves more than a passing definition, because by the end of the paper it will be carrying
a fermion's entire electroweak identity. We therefore stop and build the idea slowly, for a reader
meeting it cold.

Begin with what a Bloch sphere is not. A sphere, by itself, cannot be \emph{read}. To say that a
state points ``up'', or that a rotation went ``clockwise'', one must first lay a grid over the
sphere: choose which direction counts as $z$ (the axis the computational basis, and with it $T_3$,
refers to), choose an equatorial pair $x$ and $y$, and --- the part that turns out to matter most
--- choose the \emph{orientation} of the triple: whether $(x, y, z)$ is right-handed or
left-handed. That oriented coordinatisation is the frame. Its continuous part is a rotation's worth
of convention; its discrete kernel is a single bit per sphere, the handedness label
$\chi \in \{L, R\}$, and this bit is the Filatov datum of the previous subsection. The crucial
distinction, worth italicising: \emph{the frame is not the state}. The state is a point of the
sphere; the frame is the grid the sphere is read through --- map-grid, not territory. Two
descriptions of one state in two frames differ, and nothing physical has changed; a chart drawn
with north at the top and a chart drawn with north at the left describe the same sea.

Why, then, are frames physics at all? Because of the move this framework makes everywhere: a
\emph{lone} sphere's frame is pure gauge --- there is nothing for its orientation to be relative
to, no laboratory magnet to give ``up'' a meaning --- but \emph{relative} frame data between
spheres is physical, and entanglement legislates it. Filatov's theorem
\autocite{filatov2024towards} is exactly such legislation: an entangled pair of qubits
\emph{forces} the two spheres' frames to be opposite-handed. A frame is thus a convention per
sphere and a constraint per pair. And when three pairwise constraints meet on a triangle, they
cannot all be satisfied --- $K_3$ is not two-colourable --- so the register is frustrated, and the
inequivalent ways of resolving the frustration are \emph{discrete} invariants of the constraint
graph. Discreteness is the promotion: one cannot smooth a resolution class away by
re-conventioning, any more than one can re-convention a left glove into a right one. Continuous
choices are bookkeeping; the classes they fall into are structure.

Now count what is written on the frame, and notice that it is everything electroweak. Generation:
\emph{which} vertex of the triangle is the odd one out --- which algebra rung's frame disagrees
with the other two (\S\ref{sec:generations}). Chirality: the handedness content itself. Weak
isospin: the frustrated pair has exactly two resolutions, and $T_3 = \pm\tfrac12$ is nothing but
\emph{which resolution the channel instantiates} --- the doublet is two frame-labellings of one
frustrated structure, connected by an $SU(2)$ rotation and by nothing smaller. Colour and
hypercharge live elsewhere (on the fibres of \S\ref{sec:charges}); generation, chirality and
isospin live here. A fermion's weak identity is frame-borne.

The point sharpens into an asymmetry that the rest of the paper depends on. For a \emph{lepton},
the frame structure is certified by the state itself: the W-class register has every pair
entangled, so Filatov's theorem derives the handedness relations from the state, and the state
even co-displays the distinguished axis --- the single-qubit marginal has the non-degenerate
spectrum $(\tfrac23, \tfrac13)$ that singles out $T_3$ (Proposition~\ref{prop:w-doublet}). For a
\emph{quark}, both certificates are absent: the GHZ register has no entangled pairs, and its
single-qubit marginal is exactly $\tfrac12 I$, displaying no preferred axis whatsoever. A quark's
electroweak identity is therefore written \emph{nowhere in its local state}: it lives entirely in
the frame, as class-specification data. This is why \S\ref{sec:charges} routes the quark doublet
through the resolution classes rather than through any marginal, and it is why the extension of
Filatov's certification from entangled pairs to the GHZ class is precisely the companion clause of
Open Problem~\ref{op:filatov-extension}: the certification gap and the open problem are the same
thing, seen twice.

The asymmetry answers a question a careful reader will have been storing up: if quarks have no
entangled pairs, what does the weak force grab? The frame --- because there is nothing else on
offer, and because frames are exactly what local operations rotate. The charged current's action on
a fermion is a local unitary on the $B$ qubit, turning the register from one resolution of the
frustrated structure to the other; a turn of the grid, costing no entanglement budget and
preserving the species class exactly. Contrast the mode the vertex does \emph{not} have: attaching
to a quark by state-level entanglement, a pairwise tangle lodging on its qubits. For a saturated
register that mode is not merely absent from the dictionary --- it is forbidden by theorem, the
external pairwise capacity of a GHZ register being exactly zero (\S\ref{sec:budget-converse}).
The weak vertex on a quark has one legal mode: act on the grid, never on the territory. \emph{A
quark interacts weakly by having its frame turned, never by being held.} The worked contrast ---
the dictionary's implementation preserving the quark exactly while a pair-forming implementation
destroys it at first order --- is exhibited at the beta-decay circuit of \S\ref{sec:interactions}.

One consonance is too striking to leave unremarked, though we flag it as consonance and not yet
derivation. Empirically, the weak force is \emph{the} parity-violating interaction --- the one
force in nature that distinguishes left from right, the fact every textbook leads with. Here that
is not an anomaly to be accommodated; it is the job description. The weak force is, by
construction, the force whose entire mechanism is the rotation of oriented frames, and a force that
acts on orientation could not fail to see orientation: parity sensitivity is what ``acting on
frames'' \emph{means}. The derivation of the full left-only coupling remains the business of
Open Problem~\ref{op:chirality}; what the frame picture supplies is the reason such a force was
always going to be the odd one out under mirrors.

Finally, the geometric home, recorded as structure. The handedness labelling is an orientation
section of the product of the spheres' frame bundles; the pairwise demand is a cocycle condition on
the edges of $K_3$; the $W$ is a fibre rotation; and the odd cycle's $\mathbb{Z}_2$ holonomy is
the same object the next subsection reads as the $4\pi$ period of a spinor. That formulation is
where the Filatov extension would naturally live, and it travels with
Open Problem~\ref{op:filatov-extension}. \emph{QC view:} the frame is the per-wire basis choice
plus one handedness bit, physical only in comparisons, constrained by entanglement, classified by
the triangle. \emph{GT view:} an orientation datum on a frame bundle whose frustrated cocycle
carries generation, chirality and isospin --- the electroweak charges as topology of the grid, not
geography of the point.

\subsection{Why three spheres make a fermion}

The two-sphere representation also gives a structural reading of spin, and it must be earned
rather than asserted, because the obvious objection is strong: the Bloch sphere's most familiar
role is as the state space of a spin-$\tfrac12$ system, so a dictionary that assigns one sphere to
spin $0$ and needs three spheres to reach spin $\tfrac12$ looks inverted. The objection dissolves
once one remembers what kind of framework this is. A Bloch sphere is a spin-$\tfrac12$ state space
\emph{relative to an external reference frame} --- the laboratory magnet that gives ``up'' its
meaning. The register grants no such frame: its spheres carry internal charges, the framework is
relational throughout (\S\ref{sec:relationalism}), and the orientation of a lone sphere is
therefore unobservable --- pure gauge. Spin, here, is the transformation behaviour of what a
channel actually possesses: its \emph{relational} content, the orientations of its spheres
relative to one another. The dictionary counts exactly that, and nothing else.

Read the three classes by their constraint edges. The product register has no edge: no relative
orientation exists, nothing transforms under a reorientation, and the channel is a scalar ---
spin $0$ --- in exact consonance with the Higgs ontology of \S\ref{sec:boson-masses}, the
product-state scale-setter, pure magnitude with nothing to orient. The Bell register has one edge,
and the edge is satisfiable --- two-colourable --- so no sign lurks in it; its entire observable
content is a single relative orientation between two spheres, and a relative orientation is
vector-valued: infinitesimally an axis and an angle, the adjoint. Equivalently, by the map--state
duality of \S\ref{sec:canonical-orbits}, a Bell register \emph{is} a transport map between its
two spheres, and the physical excitations of a transporter are adjoint-valued fluctuations about
the identity --- a connection's quanta. One relation, one vector index: spin $1$.

The three-sphere register is the first configuration that can close a \emph{cycle}, and when all
three pairs are entangled --- the W class, the leptons --- the cycle is odd: the pairwise
anti-alignment constraints of the $K_3$ structure cannot all be satisfied --- that is the
frustration --- so transporting an alignment choice once around the triangle returns it flipped,
and only the second circuit restores it. A closed cycle does not add another vector; it
adds a \emph{sign}. Holonomy $-1$ under one circuit and $+1$ under two is precisely the signature
of a projective, double-valued representation --- the $4\pi$ period that defines a spinor --- and
the canonical-orbit minimality of \S\ref{sec:canonical-orbits} selects the least representation
consistent with that holonomy: spin $\tfrac12$. (Minimality chooses among representations carrying
the same holonomy; it cannot empty a channel of the relation it is, which is why the Bell class is
not spin $0$ by the same principle.) The GHZ class --- the quarks --- has no
pairwise-entangled edges, so it does not close this triangle directly; its spinor character and
doublet structure ride the same $K_3$ resolution machinery by the route \S\ref{sec:generations}
takes --- exactly as its $T_3$ already does in \S\ref{sec:charges} --- and that extension is owed
there, not assumed here. The Filatov constraint on jointly constrained spheres realises the double
cover concretely for the entangled pairs \autocite{filatov2024towards}, and the same odd-cycle
structure is what the entanglement budget renders monogamous --- the unshareability that shadows
exclusion --- so the fermion/boson split and the frustrated/unfrustrated split are, structurally,
one split seen twice: established on the W side, owed on the GHZ side.

\emph{QC view:} count the constraint edges and cycles of the register. \emph{GT view:} read off
the spin --- no relation, scalar; one relation, vector; one odd cycle, spinor. The ledger is owed
exactly two identifications, and we name them rather than gesture at them: that a $2\pi$ rotation
of the channel in the emergent space of \S\ref{sec:space-of-computations} acts on the register as
one circuit of its constraint cycle --- the bridge from internal holonomy to spatial spin --- and
that the pairwise Filatov constraint extends to three mutually constrained spheres, the assumption
disclosed above on which the $K_3$ derivation itself rests. Given the relational premise and those
two, the ladder above is derived; both are structural here, and their formal versions travel with
the generation proof and the Filatov constraint in \S\ref{sec:generations}.

\subsection{Status, and the standard objections}

The stage is an assembly, and two standard objections should be met head-on.

First, is this genuinely three qubits, or one composite state over a higher-dimensional algebra
dressed up as three? The position is that the tensor structure is a genuine tripartite product:
each of $A$, $B$, $C$ is independently a two-state system, and Proposition~\ref{prop:iterated-hopf}
is the precise sense in which the three are nested yet distinct rather than collapsed into a single
large space.

Second, does using $\mathbb{H}$ and $\mathbb{O}$ as internal algebras inherit the composite-system
difficulties of quaternionic and octonionic quantum mechanics? That objection was met in full in
\S\ref{sec:division-algebras}: the base field remains $\mathbb{C}$ throughout, the higher
algebras are internal to individual register qubits, and the $C$-qubit is handled geometrically as
$S^{15}$ precisely because $\mathbb{O}^{2}$ is not a Hilbert space. The register instantiates that
stance; nothing here reopens the question.

With the register in place and each qubit tied to its gauge sector, the next stage reads the
electric charge, weak isospin, hypercharge and colour directly off the state, recovering
$Q = T_3 + \tfrac12 Y$ as a statement about the register rather than a tabulated assignment
(\S\ref{sec:charges}).
